\section{Reading touch after a change of hardware}
\label{sec:method}
\label{sec:tracker}

\paragraph{Setting.}
A classifier is trained on labelled contacts recorded on some hardware configurations. Then the unit or the gel is replaced, and the robot keeps working. It has no labels for its new contacts, no calibration object and, in most archived data, no frame of the bare gel. What it does have is the stream of its own contacts with whatever it touches, all recorded under the new state. The method uses that stream twice: to read every test image relative to an estimate of the state (\cref{alg:tracker}), and as the target domain when the encoder is retrained.

\paragraph{The state in an image.}
We write a contact image at frame position $p$, taken on hardware configuration $h$, as $x(p) = F_h(p) + r(p)$: the state of that configuration plus the touch. \Cref{sec:signature} showed that $F_h$ is shared by every object imaged on that hardware and is replaced with the unit or the gel. It has two parts, a fine pattern and a smooth illumination field. Because each object in a typical dataset is recorded in its own session, and therefore on one configuration, a classifier can use $F_h$ as a feature of the classes recorded under it. After a change of hardware, every image then carries the feature of the wrong classes, as the injected tie on SITR shows (\cref{fig:confound}).

\begin{algorithm}[t]
\caption{Reading the robot's stream relative to its hardware state}
\label{alg:tracker}
\begin{algorithmic}[1]
\Require training mean $G$; Gaussian low-pass $L_\sigma$; change detector with thresholds from the training data
\State $\mathcal S \gets \emptyset$ \Comment{mean contact frames of the sessions since the last change}
\For{each session of the stream (one object, no label)}
  \State $\hat F \gets \big(G + \sum_{M\in\mathcal S} M\big) / \big(1 + |\mathcal S|\big)$
  \For{each episode of the session}
    \State \textbf{read} every contact crop: $x'(p) \gets x(p) - L_\sigma\big(\hat F - G\big)(p)$ \Comment{to the encoder}
    \If{the detector declares a change of hardware}
      \State $\mathcal S \gets \emptyset$; \ $\hat F \gets G$
    \EndIf
  \EndFor
  \State add the session's mean contact frame to $\mathcal S$
\EndFor
\end{algorithmic}
\end{algorithm}

\paragraph{Reading relative to the state.}
The robot estimates the current state as the mean of the contact frames it has recorded since the last change of hardware. Each earlier session contributes its mean contact frame, so every object counts once. The training mean $G$, the average of the training sessions' mean contact frames, counts as one more session. The estimate is causal. A session is read with the estimate built from the sessions before it, so the first session after a change is read relative to $G$ alone, which leaves it unchanged. Every object imprints the gel differently, so averaging over objects leaves the hardware. This is the retrospective flat-field of microscopy~\cite{singh2014illumination,peng2017basic}, taken from a robot's stream.

The reading subtracts only the smooth part of the difference between the estimate and the training mean, a Gaussian low-pass with $\sigma = 33$\,px (0.37\,mm). There are two reasons. An estimate built from a few objects still contains their average imprint, and that imprint lies mostly in the fine band. The smooth illumination field, on the other hand, is estimated well from a few objects, and it carries most of the difference between hardware states: our two units differ by \LpDistSoneStwo\ grey levels in it. The reading needs no labels, no training and no bare-gel frame. It changes the input of a trained classifier, so it can be added to any classifier, including one trained with a published adaptation method (\cref{sec:baselines}). Only the test images are read. The training images and the stream used for alignment enter as recorded (row h of \cref{tab:ablation} reads the training images as well).

\paragraph{Dating the changes.}
The estimate must restart when the hardware changes, and a robot is not always told when that happens. The detector therefore watches whole frames, which the robot records anyway. It compares the mean frame of the last $K\!=\!8$ episodes with that of the $W\!=\!30$ episodes before them in two statistics. The first is the correlation of their fine bands (190--370\,\si{\micro\metre}, at least 60\,px from any marker), which drops when the cartridge's pattern changes. The second is the mean absolute difference of their illumination fields (Gaussian, $\sigma\!=\!33$\,px). Either alone raises false alarms, because a strongly textured material lowers the fine-band correlation on any hardware and the illumination moves a little between sessions. A change is declared only when both leave the same-hardware range for $m\!=\!5$ consecutive episodes. The thresholds are the 10th percentile of the correlation and the 95th percentile of the illumination distance on the first training configuration. \Cref{sec:trackerexp} evaluates the detector.

\paragraph{Aligning the encoder with the robot's own touches.}
The reading removes an additive, smooth state. The fine pattern remains, and two sensors can also differ in ways that are not additive, such as the part of the gel the camera sees or the contrast of the image (\cref{sec:second}). For these, the encoder itself must stop separating the old hardware from the new. When the encoder is retrained, we give it the unlabelled stream from the new hardware as the target domain of an adversarial alignment. We use CDAN+E~\cite{long2018cdan} without changes to its loss. With the reading, it gave the highest accuracy of the aligners we tried, although DANN, SDAT and MIC come within the interval (\cref{sec:baselines}). A domain discriminator receives the outer product of the encoder's features and the classifier's detached predictions. When this product has more than 4096 dimensions, it is randomly projected to 1024. Each sample is weighted by $1 + e^{-H(\hat y)}$, so that confident predictions count more, and the encoder is trained against the discriminator through gradient reversal with a warm-start schedule. The objective is the cross-entropy on the labelled training crops plus the discriminator's weighted loss, with weight one, as in the Transfer-Learning-Library implementation.
